% Reconstructed from the compiled lecture-note PDF.
\chapter{The area functional and the Reinhardt control problem}
\section{Green’s theorem}

Let γ(t) = (x(t), y(t)) be a positively oriented simple closed curve. Green’s theorem gives


\begin{displaytext}
1 I 1 Z \
area(γ) = (x dy −y dx) = det(γ, γ′) dt. \
2 γ 2
\end{displaytext}

For a curve generated by g(t), write


\begin{displaytext}
γ(t) = g(t)v, g′ = gX.
\end{displaytext}

Then γ′ = gXv,


\begin{displaytext}
and, because det g = 1, \
det(γ, γ′) = det(v, Xv).
\end{displaytext}

The identity det(v, Xv) = −vTJXv


turns area into a matrix expression.


\section{Adding three synchronized arcs}

One sixth of the boundary is represented by three even-indexed curves


\begin{displaytext}
σ0(t) = g(t)s∗0, σ2(t) = g(t)s∗2, σ4(t) = g(t)s∗4,
\end{displaytext}

\begin{displaytext}
for 0 ≤t ≤tf. Their centrally symmetric copies complete the boundary. \
A direct matrix identity for any Y ∈sl2(R) is
\end{displaytext}

3 JY + (R2)TJY R2 + (R4)TJY R4 = tr(JY )I. 2


Applying Green’s theorem to the three arcs gives


Z tf area(K) = −3 tr(JX(t)) dt. 2 0


Since the star inequalities imply tr(JX) < 0, the integrand is positive.


\begin{statementbox}{Theorem}
Theorem 6.1 (Area formula). For a normalized balanced disk in circle representation,
\end{statementbox}

Z tf area(K) = −3 ⟨J, X⟩dt. 2 0


\begin{statementbox}{Example}
Example 6.2 (The unit disk). For the unit disk,
\end{statementbox}

\begin{displaytext}
π \
g(t) = etJ, X(t) = J, tf = . \
3
\end{displaytext}

\begin{displaytext}
Because J2 = −I and tr(J2) = −2,
\end{displaytext}

\begin{displaytext}
Z π/3 area(K) = −3 (−2) dt = π. \
2 0
\end{displaytext}

\section{The complete control problem}

We can now state the geometric problem as an explicit optimal-control problem.


\begin{statementbox}{Theorem}
Theorem 6.3 (Reinhardt control problem). Find a measurable control u : [0, tf] →UT , a terminal time tf > 0, and state curves
\end{statementbox}

\begin{displaytext}
g : [0, tf] →SL2(R), X : [0, tf] →sl2(R)
\end{displaytext}

that minimize −3 Z tf ⟨J, X⟩dt 2 0 subject to


\begin{displaytext}
g′ = gX, (6.1) \
[Zu, X] X′ = ⟨Zu, X⟩, (6.2) \
det X = 1, ρj(X) > 0, (6.3) \
g(0) = I, g(tf) = R, (6.4) \
X(tf) = R−1X(0)R. (6.5)
\end{displaytext}

The terminal time is free.


The state space is the tangent bundle of SL2(R), written concretely as


SL2(R) × sl2(R).


No abstract tangent-bundle machinery is needed to solve the equations: the pair (g, X) is simply a matrix and its body velocity.


\section{Why the free terminal time matters}

\begin{displaytext}
The parameter t was chosen by det X = 1, not by prescribing how long one sixth of the boundary \
must take. Different shapes can have different tf. Thus the endpoint g(tf) = R must be reached at \
a time chosen by the optimizer.
\end{displaytext}

For free-terminal-time problems, the Pontryagin Hamiltonian will vanish along every extremal. This additional equation is used repeatedly in the singular-locus and blowup analysis.


\section{Existence of an optimizer}

The original geometric problem has a minimizer by Reinhardt’s compactness argument. The control formulation also fits a standard existence theorem due to Filippov, once the relevant state region is made compact.


A simplified version of Filippov’s theorem says:


\begin{insightbox}{Standard theorem used as a black box}
\end{insightbox}

If the control set is compact, the velocity set is convex at each state, and the relevant vector fields are confined to a compact region, then the attainable sets are compact and an optimal measurable control exists.


The control simplex UT is compact. The velocity set for the upper-half-plane version of the X equation is the convex hull of the three vertex velocities. The only subtle point is that the star domain is noncompact. The paper proves that a global minimizer never approaches its ideal boundary.


\section{Why the boundary of the star domain is nonoptimal}

A state X satisfying det X = 1 and the star inequalities determines a critical hexagon with fixed midpoints s∗j. As one star inequality approaches equality, the hexagon degenerates toward a parallelogram. A convex disk whose critical hexagon is nearly a parallelogram has packing density close to 1, far above the density of the smoothed octagon. Such a state cannot belong to a global minimizer. The source makes this quantitative. It constructs a compact subset H∗∗of the star domain by deleting six neighborhoods of the three ideal vertices and three boundary sides. If a trajectory leaves H∗∗, a geometric area estimate implies


\begin{displaytext}
δ(K) > δoct.
\end{displaytext}

Therefore every global minimizer remains in H∗∗.


\begin{insightbox}{Relation to the source paper}
\end{insightbox}

This is the compactification theorem of source Chapter 5. The explicit constants, such as the horocycle height 4.5 and slope 15, are chosen for a convenient rigorous exclusion, not because


they are optimal.


\section{A useful time bound}

In upper-half-plane coordinates, developed in the next chapter, the cost integrand is at least 3. Since a minimizer has area less than the disk’s area π, one obtains


π tf < . 3


This uniform time bound, together with compactness of the X-state, controls g by Gronwall’s inequality.


\section{The conceptual achievement of the reduction}

The original unknown is a convex boundary curve. In the control problem the unknowns are:


• six scalar functions: three coordinates of g and three of X, subject to algebraic constraints; • one measurable control u(t) in a fixed triangle; • one scalar terminal time tf.


The equations are explicit rational matrix ODEs. The price is that the control can be discontinuous and may switch infinitely often. This is exactly where optimal-control theory becomes necessary.


\begin{insightbox}{Proof roadmap}
\end{insightbox}

From now on the geometric theorem will follow from a statement about controls:


global minimizer is bang-bang with finitely many switches


Indeed, each constant vertex-control segment produces straight lines and a hyperbola, and finitely many segments produce a finite smoothed polygon.


\begin{insightbox}{Chapter summary}
\end{insightbox}

Green’s theorem turns area into the linear cost −3 R ⟨J, X⟩. Together with the curvature 2 state equation and rotational endpoint conditions, this gives a free-time optimal-control problem on explicit 2 × 2 matrices. A geometric compactification excludes states near the degenerating boundary, ensuring that an optimal control exists in a compact region.

